\subsection{Decimals and Fractions}
A fraction can also be expressed as a \makeblank.
\begin{theorem}
Every fraction can be expressed as either a terminating or repeating decimal. Every terminating and repeating decimal can be expressed as a fraction.
\end{theorem}
\begin{proof}
To write a fraction $\frac{n}{d}$ as a decimal, perform the division $n\div d$. In the process of long division, we will always reach a point where the remainder is either zero, or repeats itself.

To write a terminating decimal as a fraction, write the decimal as $\frac{w}{10^n}$, where $w$ is the number obtained by deleting the decimal point and $n$ is the number of digits after the decimal point.

To write a repeating decimal as a fraction, use the algebra process demonstrated below.
\end{proof}
\begin{example}
Write each fraction as a decimal.
\begin{itemize}
\item $\frac{3}{5}=\makeblanksmall$
\item $\frac{5}{11}=\makeblanksmall$
\item $\frac{7}{12}=\makeblanksmall$
\end{itemize}
\end{example}
\begin{example}
Write each decimal as a fraction.
\begin{itemize}
\item $2.25=\makeblank[1in]$
\item $0.3\overline{18}=\makeblank[1in]$
\end{itemize}
\begin{lpic}{}{10}
\regtextleft{0}{-1}{Let $x=$};
\regtextleft{1}{-2}{digit before repeating, so multiply by};
\regtextleft{1}{-4}{digits repeat, so multiply by};
\regtextleft{0}{-7}{Subtract to eliminate the repeating part};
\end{lpic}
\end{example}